
\section{The State}




\begin{frame}{The State}
    \showmytikz[1,2,3]{../shared/tikz/tikz-flowmodel-own}
    \mycomment{
        So far, we have been looking at markets, thus at the results of exchange between households (and companies).

        The main result is the First Theorem of Welfare Economics: 

        \begin{quote}
            Under ideal conditions---perfectly competitive markets, complete information, and no externalities---any competitive equilibrium leads to a Pareto efficient allocation of resources.
        \end{quote}

        For this, the state would have to do is guarantee the workings of free markets by guaranteeing property rights, that contracts are fulfilled, that that markets are competitive (no monopolies), and that there are no externalities.

    }
\end{frame}

\begin{frame}{The State}
    \showmytikz[1,2,3,4]{../shared/tikz/tikz-flowmodel-own}
    \mycomment{
        So the state is indeed needed (and its existence justified) not only to guarantee free markets, but also to fight externalities. On top, society may well have other goals than Pareto efficiency. Fairness and equality (of chances) would be on top of the list.

        The state has many ways to influence the economy:

        \begin{itemize}
            \item securing property rights
            \item taxes
            \item regulation, such as setting standards
            \item fiscal policy, that is: expenditure
            \item monetary policy, as set by the central bank
        \end{itemize}
       
    }
\end{frame}

\begin{frame}{Levels and Institutions of the State}
    \mytext{
        \begin{itemize}
            \item Government
            \begin{itemize}
                \item Federal level
                \item State/Länder level
                \item Communal level
                \item Social Security
                \item Governmental Corporations
                \item Supranational level (EU)
            \end{itemize}
            \item Central Bank
        \end{itemize}
    }
    \mycomment{
    A typical modern state has three branches: the executive, judicative and legislative branch. On top, there are various levels of government, particularly in federal states and in the European Union.
    
    In counties with Statutory Law (as opposed to Common Law), private sector entities act under Private Law (\textit{Privatrecht}), while public sector entities fall under Public law (\textit{Öffentliches Recht}).

    Similarly, consumers may legally be treated differently than companies (e.g. \textit{BGB, HGB}) and corporations.
    }
\end{frame}

\begin{frame}{Shell Model}
    \mycomment{
    The state is not just the central government, there are many different levels of government. In Germany, most centrally collected taxes are shared 50:50 between federal and state governments. However, central and regional government may have different taxation types at their disposal. In America, excess taxes are typically levied at state, not federal, level.

    The `shell model' was developed in the framework of the European System of Accounts (ESA; German: \textit{ESVG}), so that budgets can be compared. 

    Governments may move parts of the budget into extra budget for cosmetic or practical reasons. \cite{destatisbegriffe2019}

    Governments can levy taxes, and - in general - hold a monopoly on the use of force.
    }
    \showmytikz{../shared/tikz/tikz-shellmodel.tex}
\end{frame}

\begin{frame}{Government Expenditure Ratio}
    \mytext{
            The state is a powerful actor. Just how large the role of the state in a country is can be measured by the state quota, that is the government spending ratio:
        \begin{equation}
            SQ=\frac{G}{Y}\%
        \end{equation}
        \mytable{
        \begin{tabular}{lr}
            \toprule
            Ukraine & 66\%\\
                France & 58\% \\
             Germany & 49\%  \\
             Japan & 44\%\\
             USA &  36\% \\
             Saudi Arabia & 28\%\\
             Pakistan &  20\% \\
             \bottomrule
        \end{tabular}
        }
    }
    \mycomment{The government spending ratio varies widely, with low income countries typically lower in the table. The highest ratio are offshore havens and other outliers. }
\end{frame}

% \begin{frame}{The State}
%     \mycomment{
%     The state is the institution of institutions. Only the state has the right to use force. The state is bound to the rule of law. It can raise taxes and set rules.
%     }
%     \mytext{
%     \begin{itemize}
%         \item different actor
%         \item ability to set rules (rule of law)
%         \item ability to tax
%         \item ability to issue / define money
%         \item ability to transfer
%         \item economic activity
%         \item services, education, health care, defence, infrastructure...
%         \item different levels of government
%     \end{itemize}    
%     }
%     \mycomment{Apart from the central bank, the state as the sovereign has a vast amount of power. It also serves a guarantor for the citizens rights to own property, to trade, and so forth (John \noun{Locke}). Many of these actions are of macroeconomic relevance. In microeconomics, we only study policies targeting single markets.    
%     }
    
% \end{frame}




\begin{frame}{Types of Policy}
    \mycomment{
    The state is the institution of institutions. Only the state has the right to use force. The state is bound to the rule of law. It can raise taxes and set rules.
    }
    \mytext{
    \begin{itemize}
        \item Fiscal Policy
        \begin{itemize}
            \item Government spending
            \item Taxation, tariffs
            \item Borrowing
        \end{itemize}
        \item Regulatory Policy
        \begin{itemize}
            \item restrictions
            \item obligations
            \item standards
        \end{itemize}
        \item Monetary Policy
        \begin{itemize}
            \item interest rates
            \item open market policy
        \end{itemize}
    \end{itemize}
    }
\end{frame}
 


\begin{frame}{Pareto Optimum, Pareto Criterium}
    \mytext{
    A social situation is a Pareto optimum, if nobody can be made better off without anyone else being made worse off.
    \begin{itemize}
        \item exchange leads to Pareto optimum
        \item perfect markets lead to Pareto optimum
    \end{itemize}
    Pareto criterium:
    \begin{itemize}
        \item if a situation is not Pareto-optimal, but a Pareto optimum can be achieved, act so that the optimum gets achieved
    \end{itemize}
    }
    \mycomment{
    The Pareto criterium acts as a minimalist criterium. All people will always prefer situations that makes everyone better off.

    The exchange can also be non-materialistic; I might just be better off a little poorer in monetary terms and not having to watch my neighbour suffer. In this case, a voluntary donation moves us to Pareto optimum.
    }
\end{frame}

\begin{frame}{Kaldor-Hicks criterium}
    \mytext{
    Improving some $A$ at the cost of $B$ (violating the Pareto criterium) can be justified when some of the gains of $A$ can be redistributed (later) to $B$.

    \begin{itemize}
        \item redistribution by tax and transfer
        \item redistribution by growth and future gains (of $B$)
    \end{itemize}
    }
    \mycomment{
    The Kaldor-Hicks criterium is similar to the Pareto criterium, but it requires a mechanism for the redistribution. 
    }
\end{frame}

\begin{frame}{Pareto Optimum}
    \showmytikz{../shared/tikz/tikz-pareto.tex}
    \mycomment{
    A \textbf{Pareto optimum} is the situation where no improvement to one party can be made without leaving the other party worse off.
    
    In this graph the utility of of a good allocated between two parties are plotted as $X,Y$ coordinates. 

    All points on the graph are a Pareto optimum, as there is no other point for which both $A$ and $B$ would be better off. This line is called \textbf{Pareto frontier.}

    The \textbf{Kaldor-Hicks} criterium states that even a non-Pareto-optimal situation may be desirable if and only if the disadvantage for the losing party could be compensated in another transaction. For example, a large win for $A$ and a small loss for $B$ could be compensated by a transfer, e.g. by taxing $A$ and benefitting $B$.
    }
\end{frame}

\begin{frame}{Price Ceilings}
    \showmytikz[1]{../shared/tikz/tikz-priceceiling.tex}
    \mycomment{
        Every now and again governments impose price ceilings.

        One was Emperor Diocletian with a price ceiling on bread (dictum de pretiis rerum venalium 301 n.Chr.), imposing a death penalty on high prices.
        
        The result was famine. Why?
    }
\end{frame}

\begin{frame}{Price Ceilings}
    \showmytikz[1,2]{../shared/tikz/tikz-priceceiling.tex}
    \mycomment{
        If the price is fixed below market price, some potential sellers can not sell at that price. This may be because of their own costs. 

    }
\end{frame}

\begin{frame}{Price Ceilings}
    \showmytikz[1,2,3]{../shared/tikz/tikz-priceceiling.tex}
    \mycomment{
        At the low price, demand is higher than it would be at market price.

        The result is excess demand: the line $F-G$. Many people want to buy but do not find a seller.


    }
\end{frame}

\begin{frame}{Price Ceilings}
    \showmytikz[2,3,4]{../shared/tikz/tikz-priceceiling.tex}
    \mycomment{
        The result is a loss in welfare for both consumers and producers. 

        There is no government neither.

        It is not clear how those who do get served are chosen; it might just be the ones who stand in the queue longest.

        And finally, it is not clear how this can be enforced. Policing costs resources. People are likely to trade anyway. Price ceilings tend to create black markets.
    }
\end{frame}

\begin{frame}{Raising and Using Taxes}
    \mytext{
        a tax is money you pay to the government without any specific return
        \begin{itemize}
            \item income tax
            \item inheritance tax
            \item sales tax
            \item import tax (tariffs)
            \item value added tax (VAT)
            \item ad valorem tax: percentage of value
            \item per unit tax: tax fixed per unit
            \item consumption tax: on particular goods
            \item ...            
        \end{itemize}
    }
    \mycomment{There are multiple levels of government. In Germany, most taxes are collected by the federal government, but then split $50:50$ between the federal government and the states (Länder). Some countries like the US do not levy VAT, but may have sales or consumption taxes even at communal level.
    
    % As we shall demonstrate, it is easy to raise a price for the state, by levies and taxes; however, lowering a price is hardly a possibility, since that would require increasing the quantity. This is a fundamental dilemma of economic policy.

    
    }
\end{frame}

\begin{frame}{Raising and Using Taxes}
    \showmybitmap[https://upload.wikimedia.org/wikipedia/commons/1/1c/Limbourg_Ves1JPG.jpg]{../shared/bitmaps/Limbourg_Ves1JPG.jpg}
    \mycomment{
    By definition, the government can use the money raised by taxes in whichever way they want.

    It is fundamentally impossible to link the income and the expenditure side of a complex undertaking. All the money raised (by taxes) "smells the same".
 
    Emperor (69-79 A.D.) Vespasian made money from public toilets. 
    "pecunia non olet", money doesn't smell. 
    
    You cannot tell the origins of money. This is known as "non-affection principle": all state income is pooled to cover all (equally pooled) state expenses. Tax money cannot be "earmarked".
    }
\end{frame}

\begin{frame}{Levies and Permits}
    \mytext{
        A levy is money you pay to the government with something in return.
        \begin{itemize}
            \item fishing permit
            \item emergency services levies
            \item often used like synonyms (any payment to the government)
            \item seizure of property to pay outstanding debt
            \item patent (i.e. allowance to produce or do something exclusively)
        \end{itemize}    
    }
    \mycomment{
        The exact meaning of `taxes' vs. `levies' may differ between jurisdictions. In Germany, there is a clear distinction between taxes (`Steuern') and levies (`Abgaben'). 

        Hypothecated taxes---such as road tax---take middle ground; here, a general tax is "earmarked" to serve a particular purpose, violating the non-affectation principle.
    
        % Historically, patents were just permits, like the `patent' to run a mill (in a certain area). Still, they are exclusive rights that guarantee a monopoly or a quasi monopoly, and they may or may not cost money. The distinction to simple permissions is not always clear, and the question of the price and structure of any permit is of course open to debate.
    }
\end{frame}

\begin{frame}{Taxing a Market}
    \showmytikz[]{../shared/tikz/tikz-tax.tex}
    \mycomment{
        Without intervention, the welfare effect of the market is the sum of consumer and producer surplus:

        \begin{equation}
            W_1=CS+PS
        \end{equation}

        The government could impose taxes:
        \begin{itemize}
            \item to reduce consumption (\textbf{Pigou-tax})
            \item to generate revenue
            \item to discourage producers
            \item ...
        \end{itemize}

        And this tax could be payed by the producer, or the consumer.
    }
\end{frame}

\begin{frame}{Taxing a Market}
    \showmytikz[1]{../shared/tikz/tikz-tax.tex}
    \mycomment{

        \begin{itemize}
            \item unit taxes shifts supply chain to the left
            \item price tax rotates supply chain to the left
        \end{itemize}

        The result is a new equilibrium $E_2$. The new price $p_t$ is higher. This does not in full go to the producers, because part of it is the tax that goes to the government. 

        As the tax increases the price, the new equilibrium is at a different place: 

        $$p_T < p^*+T$$.

        Thus, the tax will not increase the market price to a full extent.

        Note and consider the role of the elasticities of supply and demand on the exact size of the effects on consumer and producer.
    }
\end{frame}

\begin{frame}{Taxing a Market}
    \showmytikz[1,2]{../shared/tikz/tikz-tax.tex}
    \mycomment{
        The tax has reduced both consumer and producer surplus. This is a loss in welfare. However, the government has also generated a revenue ($GR$). 

        \begin{equation}
            W_2=CS+PS+GR
        \end{equation}

        And there is another effect: the new sum $W_2$ is smaller than the old sum $W_1$.

        The tax also creates a \textbf{dead weight loss} ($DWL$, a triangle). This is a reduction of welfare, and it is completely lost.

        It follows that \textit{per se} a non-taxed situation generates more welfare. 
    }
\end{frame}

\begin{frame}{Taxing a Market}
    \showmytikz[1,3]{../shared/tikz/tikz-tax.tex}
    \mycomment{
        If the government taxes the consumer rather than the supplier, the effect will be exactly the same.

    }
\end{frame}

\begin{frame}{Taxing a Market}
    \showmytikz[1,2,3,4]{../shared/tikz/tikz-tax.tex}
    \mycomment{
        It follows that it is \textbf{irrelevant} whether on paper the seller or the buyer has to pay the tax. In both cases, the price rises and the quantity goes down, there is a dead-weight loss and a part of the welfare now is government revenue. 

        Thus, discussions about whether or not producers `pass through' a tax (say, a value added tax) are vain, as the burden of a tax - the \textbf{tax incidence} - depends on the elasticities only.

    }
\end{frame}

\begin{frame}{Taxing a Market}
    \mycomment{
    Let us now look at another case, where demand is more inelastic and supply is highly elastic to price. 

    Generally, in this case there is a relatively high consumer surplus: the consumers love this product, no matter what is its price. Or maybe, they need it, no matter what they have to pay. Energy and food come to mind.

    if producers are plenty, and production costs are low, there will be relatively elastic supply. Again, this may be the case for oil which is cheap to produce (in many cases), and of which there is a large and very liquid market (no pun intended).
    }
    \showmytikz[1]{../shared/tikz/tikz-tax-inelastic}
\end{frame}

\begin{frame}{Taxing a Market}
    \mycomment{
    If demand is highly inelastic - and if supply is highly elastic - the burden of the tax is chiefly on the consumer (who is not willing or able to avoid it by much). Also, the tax does not disturb the volumes on the market much.

    The less elastic party (here: consumer) bears the larger part of the tax!

    This is a reason why taxes on energy and petrol are popular. The state earns much without distorting the market much. 

    This also means that a reduction of this tax will be `passed on' to the consumer to a good amount. The reason is that the consumer was bearing the taxt to begin with!
    }
    \showmytikz[1,2]{../shared/tikz/tikz-tax-inelastic}
\end{frame}

\begin{frame}{Taxing a Market}
    \mytext{
    In addition to the dead weight loss, there are other concerns:
    \begin{itemize}
        \item excess burden:
        \begin{itemize}
            \item collecting the tax is expensive
            \item redistribution is expensive
        \end{itemize}
        \item crowding out: 
        \item[] government displaces private investment
    \end{itemize}
    However, there are pro arguments:
    \begin{itemize}
        \item state corrects market imperfections
        \item state knows better who is in need
        \item[] (argument against relying on philanthropy)
        \item common goals need funding (defence, roads, education)
    \end{itemize}
    }
\end{frame}




\begin{frame}{Subsidising a Market}
    \mycomment{
    Subsidies are something like a negative tax: rather than increasing the cost, the government reduces them.
    }
    \showmytikz[1,2]{../shared/tikz/tikz-subsidy}
\end{frame}

\begin{frame}{Subsidising a Market}
    \mycomment{
    The effect is growing welfare in that market, since the price goes down and volumes go up. Both consumers and producers profit.

    The market is distorted and hence less efficient: 
    
    Producers are less interested in producing cost effective. 
    
    Consumers are less incentivised to look for substitutes. 
    }
    \showmytikz[1,2,3]{../shared/tikz/tikz-subsidy}
\end{frame}

\begin{frame}{Subsidising a Market}
    \mycomment{
    The subsidy is larger than the gain of welfare. There is dead weight loss and excess burden. 

    Also, the subsidy does not fall from the sky. The government expenditure must be met by government revenue elsewhere. There is thus a transfer from somewhere into this market. 

    If the source is income tax, that income tax might shift the demand curve down, since consumers will become poorer. However, this is going to affect every market, not just this market in question.

    Analogue to the tax case can be shown that the subsidy can be paid to the producers, shifting the supply curve, or to the consumers, shifting the demand curve, to the same overall effect.
    }
    \showmytikz[1,2,3,4]{../shared/tikz/tikz-subsidy}
\end{frame}

% \begin{frame}{Supply Shocks: Ad-valorem-tax}
%     \showmytikz[1]{../shared/tikz/tikz-supplyshock.tex}
%     \mycomment{
%         Suppose there is less supply. This could be because of an import tariff or any for any other reason.

%         This means the supply curve shifts to the left.
%     }
% \end{frame}

% \begin{frame}{Supply Shocks: Ad-valorem-tax}
%     \showmytikz[1,2]{../shared/tikz/tikz-supplyshock.tex}
%     \mycomment{
%         The shifted supply curve has reduced the available quantities and increased the price.  
%     }
% \end{frame}

% \begin{frame}{Supply Shocks: Ad-valorem-tax}
%     \showmytikz[2,3]{../shared/tikz/tikz-supplyshock.tex}
%     \mycomment{
%          The shift hits mostly the consumer (consumer surplus). However, the new producer surplus $PS_2$ may not be larger than the old one. And welfare in total - the sum of producer and consumer surplus - is also down.
%     }
% \end{frame}

% \begin{frame}{Supply Shocks: countermeasures}
%     \showmytikz[2,3,4]{../shared/tikz/tikz-supplyshock.tex}
%     \mycomment{
%          Suppose the government reduces tax on the good (or subsidises). That would shift the supply curve downwards. 
%     }
% \end{frame}

% \begin{frame}{Supply Shocks: countermeasures}
%     \showmytikz[2,3,4,5]{../shared/tikz/tikz-supplyshock.tex}
%     \mycomment{
%          This can indeed restore the welfare status quo for the consumer - albeit at a (huge) price for the government. 

%          The yellow area is the revenue the government \textit{would have made with the tax at the new price}, which it now loses. This is smaller than the overall increase of welfare. 

%          The graph does not show the governments opportunity costs, that is the revenue it had at the old price. 
%     }
% \end{frame}

\begin{frame}{Taxation: Laffer Curve}
    \showmytikz{../shared/tikz/tikz-laffer.tex}
    \mycomment{
        Assume the government tries to maximise government revenue. (As shown is this is not the same as maximum welfare, due to dead weight loss and excess burden.)
        
        At (income) tax of zero, government revenue is zero.

        At a tax rate of 100\%, all activities will stop. So government revenue is zero. 

        So there must be a maximum somewhere in between.

    }
\end{frame}

\begin{frame}{Taxation: Laffer Curve}
    \showmytikz[1]{../shared/tikz/tikz-laffer.tex}
    \mycomment{
        The exact shape of the Laffer curve is disputed / unknown. Most importantly, nobody is certain where the maximum lies.
        
        However, economic activity will not immediately respond to neither tax increase nor decrease. 
        
        And there must be substitution, income, tax avoidance, and macroeconomic effects - plus the question of how the revenue is spent.
    }
\end{frame}

\begin{frame}{Ramsey Rule}
    \mytext{
    \begin{itemize}
        \item government can tax different goods
        \item if price elasticity is high $\rightarrow$  tax avoidance
        \item $\Rightarrow$ most taxes where elasticity is lowest
        \item a lump-sum tax (per head) cannot be avoided at all $\rightarrow$ no distortion
        \item Ramsey pricing also serves as a profit maximising rule for optimum pricing in a monopoly
    \end{itemize}
    }
    
\end{frame}

\begin{frame}{Distribution}
    \showmytikz{../shared/tikz/tikz-gini.tex}
    \mycomment{
    The income of an economy is the sum of everything produced. It is a \textbf{flow variable.}

    The wealth of an economy is the sum of all assets. It is a \textbf{stock variable.}

    Income and wealth, anything really, can be more or less equally distributed. A measure for the equality of a distribution is the Gini coefficient. At $G=0$, the society is perfectly equal, at $G=1$ it is perfectly unequal.

    At perfect equality, there is no incentive to do anything. At perfect inequality, neither.
    }
\end{frame}

\begin{frame}{Public Choice: Median Voter Theorem}
    \showmytikz{../shared/tikz/tikz-medianvoter.tex}
    \mycomment{
        The idea is to sort all voters from left to right (on whatever the issue). That could be ranking them by the answer to the question `how high shall we tax this use of a public good'.

        Assume: to reach a consensus, a simple majority is needed.
    }
\end{frame}

\begin{frame}{Public Choice: Median Voter Theorem}
    \showmytikz[1]{../shared/tikz/tikz-medianvoter.tex}
    \mycomment{
        Any majority must necessarily involve the `guy in the middle', the median voter.

        This explains why political parties are competing for the middle - while at the same time polarising to the left or right.

        One problem is the inconsistency of the ranking regareding different topics. This can lead to instable majorities (Condorcet paradoxon), and there might not be a solution (Arrow's theorem).

        Note that there are other voting systems than simple majority voting (Condorcet methods).
    }
\end{frame}

\begin{frame}{Procurement and Tenders}
    \mytext{
    \begin{itemize}
        \item call for proposals, call for bids...
        \item should you reward the highest bidder?
        \item imagine:
        \begin{itemize}
            \item bidder A offers to do it for 2 mio
            \item bidder B offers to do it for 1.8 mio
            \item bidder C offers to do it for 200.000
        \end{itemize}
        \item sometimes in pubic procurement a law forces the authorities to reward contracts to the highest bidder (cheapest provider).
    \end{itemize}
    }
    \mycomment{The bid of bidder C might not even be realistic. You end up having to deal with a bidder that is incompetent, understating the real costs, and/or might even go bankrupt in the process. The overall costs may then exceed that of the higher bidders.}
\end{frame}

\begin{frame}{Auction Theory}
    \mycomment{
    Auctions can be designed in different ways.  The most common type, the English auction, incentivises too high bids and the winner's curse: there is no surplus left for the winner. 

    The Vickrey auction is the best system, as it incentivises to reveil one's real preferences, and leaves a surplus to the winner.
    }
    \mytext{
    \begin{itemize}
        \item open ascending-bid auctions (English auction)
        \item open descending (Dutch auction)
        \item all-pay auction (common in fundraising)
        \item open bid vs. closed bid auctions
        \item first-price vs. second-price auctions
        \item winner's curse: no surplus for the winner
        \item Vickrey auction: second price sealed-bid
        \item example spectrum auctions
        \item prequalification stages etc.
    \end{itemize}
    }
\end{frame}

% \section{External Effects}

